\documentclass[10pt,a4paper]{amsart}
\usepackage[margin=19mm]{geometry}
\usepackage{amsmath,amssymb,mathtools}
\usepackage[scr=rsfs,cal=euler]{mathalfa}
\usepackage{xfrac}
\usepackage[unicode,hidelinks,bookmarks=false]{hyperref}
\newcommand{\ie}{i.e.\ }
\newcommand{\cf}{cf.\ }
\newcommand{\resp}{respectively\ }
\newcommand{\Subsection}{\subsection}
\providecommand{\nobreakdash}{\nobreak\hbox{-}}
\setlength{\emergencystretch}{2em}
\multlinegap0pt
\usepackage{amssymb}
\usepackage{xcolor}
\usepackage{mathtools}
\newtheorem{lemma}{Lemma}
\newcommand{\E}{\mathbb{E}}
\renewcommand{\O}{\Omega}
\newcommand{\R}{\mathbb{R}}
\newcommand{\be}{\beta}
\newcommand{\bul}{\bullet}
\newcommand{\dl}{\delta}
\newcommand{\dom}{\textup{dom}}
\newcommand{\fH}{\mathfrak{H}}
\newcommand{\fm}{\mathfrak{m}}
\newcommand{\ind}{\mathbf 1}
\newcommand{\noi}{\noindent}
\newcommand{\wt}{\widetilde}
\title{Central limit theorem for stochastic nonlinear wave equation with pure-jump L\'evy white noise}
\author{Raluca M. Balan, Guangqu Zheng}
\date{Source: arXiv:2504.18672v1 (CC BY 4.0)}
\begin{document}
\maketitle

\noindent\emph{Source and licence.} Central limit theorem for stochastic nonlinear wave equation with pure-jump L\'evy white noise. Version arXiv:2504.18672v1 (CC BY 4.0), distributed by arXiv under the Creative Commons Attribution 4.0 International licence, http://creativecommons.org/licenses/by/4.0/, as recorded in the arXiv licence metadata for this exact version and shown on the abstract page https://arxiv.org/abs/2504.18672v1. Full licence notice: \texttt{licenses/arxiv-2504.18672-CC-BY-4.0.txt}.\par\smallskip

\noindent\textit{Editorial scope.} Counted unit: the article's lemma lem\_Heis (source0.tex lines 2078--2104). The article's complete original proof of it is delivered with it (source0.tex lines 2107--2274, one \texttt{proof} environment of 168 lines). No mathematics is written, repaired or re-derived: the statement and the proof are the article's own lines, in the article's own notation, and no retained line is substituted or re-flowed.  Retained uncounted as the source-local context the proof needs: lines 1963--1965; lines 1973--1975; lines 1980--2018; lines 1986--1988.  Nothing was omitted from any retained span, so the package carries no omission record.  No retained line contains a citation, so the wrapper carries no bibliography and no bibliography entry.  No retained line contains a figure environment, an image-inclusion command or drawing code, so no image is delivered and the package declares no asset dependency.  Editorial formatting only, in addition: an AMS \texttt{amsart} preamble that restates the article's own declarations and macros used by the retained lines, namely the environments \texttt{lemma} on separate counters, together with the standard packages those lines need.  The retained environments print in this document as Lemma 1 in order of appearance, whereas the article prints the same items with the numbering of its own full document.  The complete native source of the article as distributed in the arXiv e-print for this exact version is bundled unchanged as \texttt{sources/177111/source0.tex}.
\begin{align}\label{227}
V(\xi) = h_0(\xi) + \sum_{n=1}^\infty I_n\big(  h_n(\xi, \bul)  \big),
\end{align}

\begin{align}\label{227b}
\E\big[ \| V\|_\fH^2 \big] = \sum_{n\geq 0} n! \| h_n \|^2_{\fH^{\otimes (n+1)}} < +\infty.
\end{align}

\begin{lemma} \label{lem_dl}
Let  $V\in L^2(\O;\fH)$ and $V(\xi)$ have the chaos expansion as in  \eqref{227}-\eqref{227b}.
Then, $V\in\dom(\dl)$
if and only if 

 \noi
\begin{align}  \label{227d}
\sum_{n=1}^\infty n! \| \wt{h}_{n-1} \|_{\fH^{\otimes n}}^2 < \infty.
\end{align}

\noi
For $V\in\dom(\dl)$, we have
 
 
 \noi
\begin{align}   \notag%  \label{227c}
  \dl(V) = \sum_{n=1}^\infty I_n( \wt{h}_{n-1} ).
\end{align}
Here, $\wt{h}_n$ is the symmetrization of $h_n$ in all the $(n+1)$ arguments:

\noi
\begin{align}  \label{def-wth}
\begin{aligned}
\wt{h}_n(\xi_1,\ldots,\xi_n,\xi_{n+1})
&=
\frac{1}{n+1}\sum_{i=1}^{n+1}h_n(\xi_i, \pmb{\xi_{\neq i} } ),
%&=\frac{1}{n+1}\sum_{i=1}^{n+1}h_n(\xi_1,\ldots,\xi_{i-1},\xi_{i+1},\ldots,\xi_n,\xi_{n+1},\xi_i)
\end{aligned}
\end{align}

\noi
where $\pmb{\xi_{\neq i} } $ stands for the $n$ arguments $(\xi_1, ... , \xi_{i-1}, \xi_{i+1}, ... ,\xi_{n+1})$
when $i\geq 2$,  and for $(\xi_2, ... ,\xi_{n+1})$ when $i=1$.\footnote{Note that the canonical 
symmetrization of $h_n$ reduces to the above form \eqref{def-wth}, since
$h_n(\xi, \xi_1, ... , \xi_n)$ is symmetric in the $n$ arguments $(\xi_1, ... , \xi_n)$
for any fixed $\xi$.}


\end{lemma}

\begin{align}  
\sum_{n=1}^\infty n! \| \wt{h}_{n-1} \|_{\fH^{\otimes n}}^2 < \infty.
\end{align}

\begin{lemma}\label{lem_Heis}\textup{(Heisenberg's commutation relation)} 
%
%
%
Let $V\in\dom(\dl)$ as in Lemma \ref{lem_dl}.
Assume $\dl(V)\in\dom(D)$ and $V(\xi_0)\in\dom(D)$
for $\fm$-almost every $\xi_0\in\mathbf{Z}$ such that 

\noi
\begin{align}\label{cond_H1}
\E\int_{\mathbf{Z}^2} \big| D_{\xi_0} V(\xi) \big|^2 \fm(d\xi) \fm(d\xi_0) <\infty.
\end{align}

\noi
 Then,     $D_{\xi_0}V  \in \dom(\dl)$ for $\fm$-almost 
every $\xi_0\in\mathbf{Z}$ with



 

\noi
\begin{align}\label{com_rel}
D_{\xi_0}  \dl(V)=  \dl(  D_{\xi_0} V) + V(\xi_0) .
\end{align}

\end{lemma}

\begin{proof}

% (i)
%Suppose  that  $V\in\dom(\dl)$ has the expression \eqref{227} 
%and that $V$ is predictable. Then, there exists a sequence 
%$V^{(k)}\in\dom(\dl)$ such that 
%$V^{(k)}\to V$ in $L^2(\O; \fH)$
%and $V^{(k)}$ has the following form
%\begin{align}  \label{228a}
%V^{(k)}(r, y, z) = \sum_{j\in \be_k} Y_j \ind_{(a_j, b_j] \times A_j \times \Gamma_j}(r, y, z), 
%\end{align}
%
%
%\noi
%where $\be_k$ is a finite index set, $a_j < b_j$, 
%$A_j\times \Gamma_j\in \mathcal{B}(\R)\times \mathcal{B}(\R_0)$, 
%and $Y_j\in\dom(D)$
%is $\mathcal{F}_{a_j}$-measurable. One can further assume $Y_j$ admits the following chaos expansion
%\begin{align}  \label{228b}
%Y_j = \sum_{n=0}^\infty I_n( g_{n,j})
%\quad
%\text{with $g_{n,j}\in \fH^{\odot n}$}.
%\end{align}
%Combining \eqref{228a} and \eqref{228b} yields,
%with  $\xi = (r,y,z)$, 
% 
% \noi
%\begin{align}  \label{228c}
%\begin{aligned} 
%&V^{(k)}(r, y, z) 
%= \sum_{n=0}^\infty I_n( h_{n,k}(\xi, \bul))\\
%&\text{with}
%\quad  
%   h_{n, k}(\xi) := \sum_{j\in \be_k} g_{n,j}(\bul) \ind_{(a_j, b_j] \times A_j \times \Gamma_j}(\xi)\in\fH^{\odot n}
%\end{aligned}
%\end{align}
%
%\noi
%It is clear that $h_{n, k}$ converges to $h_n$ in $\fH^{\otimes (n+1)}$
%as $k\to\infty$, since $V^{(k)}\to V$ in $L^2(\O; \fH)$.
%Therefore, we can deduce from Fatou's lemma that 
%
%\noi
%\begin{align*}
%\sum_{n=1}^\infty n!  \| \wt{h}_{n-1} \|^2_{\fH^{\otimes n}}
%&= \sum_{n=1}^\infty n!  \lim_{k\to\infty} \| \wt{h}_{n-1, k} \|^2_{\fH^{\otimes n}} \\
%&\leq  \liminf_{k\to\infty}   \sum_{n=1}^\infty n! \| \wt{h}_{n-1, k} \|^2_{\fH^{\otimes n}} 
%=  \liminf_{k\to\infty}  \| \dl(V^{(k)} ) \|_2^2 = \| \dl(V) \|_2^2 .
%\end{align*} 
% That is, the condition \eqref{227d} holds true, and 
% the above ``$\leq \liminf$'' can be replaced by ``$ = \lim$''. 
% 
 
  Let $V\in\dom(\dl)$  have the expression \eqref{227}.
 Then,   the assumption $V(\xi)\in \dom(D)$ 
 implies 
 
 \noi
 \begin{align}\label{DxV}
 D_{\xi_0} V(\xi) = \sum_{n=1}^\infty n I_{n-1}\big( h_n(\xi, \xi_0, \bul) \big).
 \end{align}

\noi
Since 
the condition \eqref{cond_H1}    implies

\noi
\begin{align}
\int_{\mathbf{Z}^2} \E \Big[ \big| D_{\xi_0} V(\xi) \big|^2 \Big]  \fm(d\xi) \fm(d\xi_0) 
&=  \int_{\mathbf{Z}^2}  \sum_{n=1}^\infty n^2 (n-1)! \| h_n(\xi, \xi_0, \bul) \|^2_{\fH^{\otimes (n-1)}}  
\fm(d\xi) \fm(d\xi_0) \notag   \\
&= \sum_{n=1}^\infty n  n! \| h_n \|^2_{\fH^{\otimes (n+1)}},
\label{cond_H2}
\end{align}  

\noi
we get 

\noi
\begin{align}
\sum_{n=1}^\infty n  n! \| h_n \|^2_{\fH^{\otimes (n+1)}} < \infty.
 \notag% \label{cond_H2b}
\end{align}  

Before proving  $D_{\xi_0} V\in\dom(\dl)$,
we point out that the assumption $\dl(V)\in\dom(D)$ leads to 

\noi
\begin{align}
\sum_{n=1}^\infty n n! \| \wt{h}_{n-1}\|^2_{\fH^{\otimes n}} <\infty
\quad\text{or equivalently}
\quad
\sum_{n=1}^\infty (n+1)^2 n! \| \wt{h}_{n}\|^2_{\fH^{\otimes (n+1)}} <\infty.
\label{cond_H3}
\end{align}

\noi
And it is also immediate that 

\noi
\begin{align}
D_{\xi_0} \dl(V) = \sum_{n=1}^\infty n I_{n-1}\big(  \, \wt{h}_{n-1}(\xi_0) \big), 
\label{Dxdl}
\end{align}

\noi
where $ \wt{h}_{n-1}(\xi_0) $ is obtained by first symmetrizing $h_{n-1}$
and then fixing the first argument to be $\xi_0$.

Now let us show $D_{\xi_0} V\in\dom(\dl)$.
For this, we apply Lemma \ref{lem_dl}.
We first rewrite \eqref{DxV} as follows:


\noi
\begin{align*}
 D_{\xi_0} V(\xi)  =  \sum_{n=0}^\infty I_n\big( g_{\xi_0, n}(\xi, \bul) \big)
 \quad
 {\rm with}
 \quad
 g_{\xi_0, n}(\xi, \bul) := (n+1) h_{n+1}(\xi, \xi_0, \bul).
\end{align*}
 
\noi
It is not difficult to see that the canonical symmetrization of $ g_{\xi_0, n}$ is given by 

\noi
\begin{align}\label{syms}
\begin{aligned}
\wt  g_{\xi_0, n}(\xi_{n+1},  \xi_1, ... , \xi_n) 
&= \sum_{j=1}^{n+1} h_{n+1}(\xi_j,  \pmb{\xi_{\neq j}}) \\
&= \bigg(  \sum_{j=0}^{n+1} h_{n+1}(\xi_j,  \pmb{\xi_{\neq j}})  \bigg) - h_{n+1}(\xi_0, \xi_1, ... , \xi_{n+1}) \\
&= (n+2) \wt{h}_{n+1}(\xi_0, \xi_1, ... , \xi_{n+1}) - h_{n+1}(\xi_0, \xi_1, ... , \xi_{n+1}),
\end{aligned} 
\end{align} 
 
 
 
 
 
 
 \noi
 where 
 $\pmb{\xi_{\neq j}}$ stands for the $(n+1)$ arguments 
 $(\xi_0, ...,  \xi_{j-1},\xi_{j+1}, ...,  \xi_{n+1})$.
 Therefore, it follows from \eqref{syms} that 
 
 \noi
 \begin{align*}
 \sum_{n\geq 1} n! \| \wt  g_{\xi_0, n-1} \|^2_{\fH^{\otimes n}}
 \leq 
 2  \sum_{n\geq 1} n! 
 \Big[   (n+1)^2 \|  \wt{h}_{n}(\xi_0, \bul) \|^2_{\fH^{\otimes n}}
  + \| h_{n}(\xi_0, \bul) \|^2_{\fH^{\otimes n}} \Big]   ,
 \end{align*}
 
 \noi
 which is integrable with respect to $\fm(d\xi_0)$ in view of \eqref{cond_H2} and \eqref{cond_H3}.
 Thus, $D_{\xi_0} V\in\dom(\dl)$
 for $\fm$-almost every $\xi_0\in\mathbf{Z}$ and
 \begin{align}\label{Dx0V}
 \dl\big( D_{\xi_0}V \big) = \sum_{n=1}^\infty I_n\big( \wt  g_{\xi_0, n-1}  \big).
 \end{align}
Hence the desired commutation relation $D_{\xi_0} \dl(V) = \dl( D_{\xi_0}V) + V(\xi_0)$
follows immediately 
from \eqref{Dxdl}, \eqref{Dx0V}, \eqref{227}, and \eqref{syms}. \qedhere
 
\end{proof}

\end{document}
